\subsection{Text-to-image Generation}
\label{subsec:t2i}
\label{subsec:p3_s8_1_1}

Text-to-image (T2I) generation synthesizes an image from a natural-language description.
For the rest of this chapter its importance lies less in the images themselves than in the prior it provides: a high-capacity mapping from compositional language to the visual space of objects, surfaces, layouts, and viewpoints, learned once at web scale and then reused.
The spatial controllers of §\ref{subsec:p3_s8_1_2} and the editors of §\ref{subsec:p3_s8_1_3} impose constraints on this prior rather than replacing it, and the text-to-3D pipelines of §\ref{sec:p3_s8_2} distill it through score distillation or multi-view consistency objectives; whatever relational competence the prior has, or lacks, therefore propagates downstream.
The field has passed through three architectural paradigms, from adversarial training to autoregressive token modeling to iterative denoising, with diffusion and its flow-matching generalization now setting the standard for quality and prompt fidelity~\cite{rombach2022ldm,peebles2023dit}.


\subsubsection{Task Definition}
We formulate text-to-image generation in a paradigm-agnostic form that exposes the structural choices shared by modern systems.
Let $\mathcal{X}$ denote the space of tokenized text prompts and $\mathcal{Y}=[0,1]^{H\times W\times 3}$ the pixel-space image domain.
Almost all modern T2I systems operate not directly on $\mathcal{Y}$ but on a compressed \emph{latent space} $\mathcal{Z}\subset\mathbb{R}^d$ provided by an autoencoder pair
\begin{equation}
E:\mathcal{Y}\rightarrow\mathcal{Z},\qquad D:\mathcal{Z}\rightarrow\mathcal{Y},\qquad D\!\circ\!E\approx \mathrm{id}_{\mathcal{Y}},
\label{eq:t2i_autoencoder}
\end{equation}
where $E$ encodes a training image into a low-dimensional code and $D$ decodes a code back into a pixel image.
A text encoder $\phi:\mathcal{X}\rightarrow\mathcal{E}$ maps the prompt into a conditioning embedding $\phi(\mathbf{x})\in\mathcal{E}$.
The core learning problem is then to train a conditional latent sampler
\begin{equation}
G_\theta:\mathcal{E}\rightarrow\mathcal{Z},\qquad \hat{\mathbf{z}}\sim p_\theta\!\big(\mathbf{z}\,\big|\,\phi(\mathbf{x})\big),
\label{eq:t2i_generator}
\end{equation}
so that, after decoding, the resulting image $\hat{\mathbf{y}}=D(\hat{\mathbf{z}})$ is both photorealistic and semantically aligned with the prompt $\mathbf{x}$.
The three paradigms covered below differ chiefly in how the conditional distribution $p_\theta(\mathbf{z}\mid\phi(\mathbf{x}))$ is parameterized and optimized; the autoencoder pair $(E,D)$ and the text encoder $\phi$ are largely shared design choices across paradigms.
Figure~\ref{fig:p3_s8_1_1_framework} schematizes this paradigm-agnostic pipeline and the substrate that the three paradigms share.

\begin{figure}[!ht]
\centering
\includegraphics[width=\linewidth]{figures/part3/section8/t2i_framework_and_paradigms.pdf}
\caption{The latent-generator view of text-to-image generation. Top: the shared pipeline of Eqs.~\eqref{eq:t2i_autoencoder} and~\eqref{eq:t2i_generator}; the encoder $E$ and the dashed paths are used only in training. Bottom: the sampler $G_\theta$ as (a) a GAN, (b) an autoregressive token model, and (c) diffusion or flow matching. Not every architecture uses a separate autoencoder; the corgi is a schematic illustration.}
\label{fig:p3_s8_1_1_framework}
\end{figure}


\subsubsection{Methodology}

Three generative paradigms, taken in the order in which each dominated T2I, parameterize the sampler $G_\theta$ of Eq.~\eqref{eq:t2i_generator} differently; the bottom row of Figure~\ref{fig:p3_s8_1_1_framework} draws them against the common pipeline.
A second set of advances refines the sampler without redefining it: backbone evolution, sampling acceleration, and the migration of leading systems to proprietary stacks.
Structural conditioning interfaces, the fourth such refinement, are the subject of §\ref{subsec:p3_s8_1_2} and are not repeated here.

\paragraph{Generative Adversarial Networks.}\label{para:p3_s8_1_1_gan}
GAN-based T2I models instantiate $G_\theta$ as a feed-forward generator $G_\theta(\mathbf{n},\phi(\mathbf{x}))\!\to\!\hat{\mathbf{z}}$ driven by an auxiliary noise vector $\mathbf{n}\sim\mathcal{N}(0,I)$; the conditional distribution $p_\theta(\mathbf{z}\mid\phi(\mathbf{x}))$ is defined only implicitly, as the pushforward of $\mathbf{n}$ through $G_\theta$.
Training follows the two-player game of \citet{goodfellow2014gan}, pitting the generator against a discriminator $D_\psi$,
\begin{equation}
\min_\theta \max_\psi\ \mathbb{E}_{\mathbf{y}\sim p_\text{data}}\!\big[\log D_\psi(\mathbf{y},\phi(\mathbf{x}))\big] + \mathbb{E}_{\mathbf{n}}\!\big[\log\!\big(1-D_\psi(D(G_\theta(\mathbf{n},\phi(\mathbf{x}))),\phi(\mathbf{x}))\big)\big],
\label{eq:t2i_gan_objective}
\end{equation}
in which $D_\psi$ scores both photorealism and text-image consistency.
The result is sharp high-frequency detail and single-pass inference, bought at the price of unstable training and mode collapse.
Within this framework, StackGAN~\cite{stackgan2017} generated a low-resolution image first and refined it with a second GAN conditioned on both the text and the coarse image; AttnGAN~\cite{attngan2018} let the generator attend to individual words while rendering different regions, with an auxiliary image-text matching loss (DAMSM) that reduced attribute errors such as wrongly colored objects; XMC-GAN~\cite{xmcgan2021} replaced that loss with cross-modal contrastive learning and improved multi-object prompts.
By 2021 diffusion models had overtaken these systems on COCO, chiefly because they train stably and cover more of the image distribution, and research attention moved on; the GAN emphasis on high-frequency detail survives in hybrid and adversarially fine-tuned samplers.

\paragraph{Autoregressive Models.}\label{para:p3_s8_1_1_ar}
Autoregressive (AR) models discretize the latent space: the autoencoder $E$ is realized as a vector-quantized variational autoencoder (VQ-VAE) whose output is a sequence of discrete tokens $\mathbf{z}=(z_1,\dots,z_N)\in\mathcal{V}^N$ over a learned codebook $\mathcal{V}$, and $G_\theta$ is factorized in causal raster order,
\begin{equation}
p_\theta\!\big(\mathbf{z}\,\big|\,\phi(\mathbf{x})\big) \;=\; \prod_{n=1}^{N} p_\theta\!\big(z_n\,\big|\,z_{<n},\,\phi(\mathbf{x})\big),
\label{eq:t2i_ar_factorisation}
\end{equation}
trained by maximum likelihood on the next-token cross-entropy.
T2I thus becomes sequence-to-sequence translation, and the transformer that performs it inherits the compositional competence of language models: DALL-E~\cite{dalle2021} modeled text and image tokens as a single sequence and showed zero-shot generalization across prompt types.
The cost is decoding latency, one forward pass per token, which grows with resolution.
Parti~\cite{parti2022} scaled the transformer to 20B parameters, generated at $256\times256$, and reached $1024\times1024$ with a convolutional super-resolution stage; Muse~\cite{chang2023muse} abandoned strict sequential order for masked token prediction, filling in many tokens per step over a handful of rounds, and matched diffusion quality on several benchmarks at a fraction of the sampling time.
AR generation remains attractive where images are produced jointly with text (interleaved documents, storyboards) and a single sequence model handles both.

\paragraph{Diffusion Models.}\label{para:p3_s8_1_1_diffusion}
Diffusion models realize $G_\theta$ as iterative denoising rather than a single feed-forward pass.
A fixed \emph{forward} process corrupts the clean latent $\mathbf{z}_0=E(\mathbf{y})$ along a schedule $(\alpha_t,\sigma_t)_{t=0}^{T}$, $\mathbf{z}_t=\alpha_t \mathbf{z}_0+\sigma_t\epsilon$ with $\epsilon\sim\mathcal{N}(0,I)$ and $\mathbf{z}_T\!\approx\!\mathcal{N}(0,I)$; a noise-prediction network $\epsilon_\theta$ learns to invert it under the denoising score-matching objective
\begin{equation}
\min_\theta\ \mathbb{E}_{t,\mathbf{z}_0,\epsilon}\!\left[\left\|\epsilon-\epsilon_\theta\!\big(\mathbf{z}_t,\,t,\,\phi(\mathbf{x})\big)\right\|_2^2\right],
\label{eq:t2i_ldm_objective}
\end{equation}
which is the diffusion-specific realization of the training objective implicit in Eq.~\eqref{eq:t2i_generator}.
At inference the model starts from $\mathbf{z}_T\sim\mathcal{N}(0,I)$, applies $\epsilon_\theta$ for $t=T,\dots,1$, and decodes $\hat{\mathbf{y}}=D(\hat{\mathbf{z}}_0)$; each step corrects a small error, which is why the final samples are both detailed and faithful.
Running this process in the latent space of Eq.~\eqref{eq:t2i_autoencoder} rather than in pixels, as Latent Diffusion Models~\cite{rombach2022ldm} did (\eg, $4\times64\times64$ instead of $3\times512\times512$ for a 512px image), cut memory and compute enough to train on commodity hardware and to release Stable Diffusion openly.
Conditioning the denoiser on embeddings from a pretrained language-image model, as in GLIDE~\cite{glide2022} and DALL$\cdot$E~2~\cite{dalle2_2022}, improved prompt alignment, and classifier-free guidance~\cite{ho2022classifierfree} supplied a single knob that trades diversity against fidelity by extrapolating the conditional prediction away from the unconditional one.
The remaining weakness is sampling cost, dozens of sequential network evaluations per image; higher-order ODE samplers and the distillation methods below address it.
Diffusion or flow-matching cores now underlie most leading systems, with the native multimodal autoregressive stacks discussed under proprietary systems the main exception.

\paragraph{Flow Matching and Rectified Flow.}\label{para:p3_s8_1_1_fm}
Rather than a fourth paradigm, flow matching is diffusion reformulated in continuous time.
It parameterizes $G_\theta$ as a learned \emph{velocity field} $v_\theta(\mathbf{z}_t,t,\phi(\mathbf{x}))$ instead of a noise predictor: training picks the straight interpolation $\mathbf{z}_t=(1-t)\mathbf{z}_0+t\mathbf{z}_1$ between a clean latent $\mathbf{z}_0=E(\mathbf{y})$ and a Gaussian sample $\mathbf{z}_1\!\sim\!\mathcal{N}(0,I)$ and regresses the velocity,
\begin{equation}
\min_\theta\ \mathbb{E}_{t,(\mathbf{z}_0,\mathbf{z}_1)}\!\left[\left\|v_\theta(\mathbf{z}_t,t,\phi(\mathbf{x}))-(\mathbf{z}_1-\mathbf{z}_0)\right\|_2^2\right],
\label{eq:t2i_fm_objective}
\end{equation}
and inference integrates the probability-flow ODE $\mathrm{d}\mathbf{z}_t/\mathrm{d}t=v_\theta(\mathbf{z}_t,t,\phi(\mathbf{x}))$ from $t=1$ back to $t=0$.
The formulation was introduced concurrently by~\citet{lipman2023flow} (Flow Matching) and~\citet{liu2023rectifiedflow} (Rectified Flow); \citet{albergo2023building} gave the unifying stochastic-interpolant view in which diffusion is the special case of a noisier interpolation schedule.
Three consequences matter in practice. Straight-line interpolations yield nearly straight ODE trajectories, so matched quality is reached in noticeably fewer steps (genuine one-step generation still needs a reflow or distillation stage); the objective is a single regression with no noise-schedule design and no per-schedule guidance recalibration; and the loss scales smoothly with model size and resolution.
Stable Diffusion~3~\cite{sd3rft2024} scaled the recipe to T2I on a DiT-style transformer backbone and reported the best GenEval scores among open models together with favorable human preference on Parti-prompts; the Flux family followed the same recipe.
Flow matching composes with every refinement below as readily as a noise-predictor core does and has become the default continuous-time formulation in 2024-onwards systems.

\paragraph{Backbone Evolution: From U-Net to Diffusion Transformers.}
Early diffusion models used the U-Net for its multiscale feature fusion and locality bias.
The Diffusion Transformer (DiT) of~\citet{peebles2023dit} replaced it with a stack of transformer layers over latent patch tokens, giving every denoising step a global receptive field, and showed that FID improves monotonically with depth, width, and token count along predictable scaling laws, reaching FID 2.27 on class-conditional ImageNet at $256^2$.
Attention cost grows quadratically with resolution, so hybrid CNN-transformer designs and efficient attention remain active; the DiT lineage nonetheless underlies SD3 and most 2024-onwards systems, and it points toward a single transformer that generates several modalities, differing only in tokenization and training data.

\paragraph{Sampling Acceleration and Distillation.}
Fifty or more denoising steps rule out interactive use and low-power deployment.
\citet{salimans2022progressive} distilled a $T$-step teacher into a $T/2$-step student and iterated the halving to reach $1$--$4$ steps.
Consistency Models~\cite{song2023consistency} train a network to map a noisy latent at any timestep directly to the clean image, either from scratch with a consistency loss or by distilling a pretrained diffusion model, and Latent Consistency Models~\cite{lcm2023} applied the idea to Stable Diffusion's latent space, matching $50$-step quality in $2$--$4$ steps.
All of these are post-training procedures that leave the architecture untouched, so any strong backbone can be made fast after the fact.

\paragraph{Proprietary and Commercial Systems.}
Leading systems have migrated to proprietary, industrially scaled stacks of three kinds.
Native multimodal models such as OpenAI's GPT~Image~1.5~\cite{openai2025images} and Google's Gemini~3~Pro~Image (Nano~Banana~Pro) run a unified transformer over interleaved text and image tokens and emphasize semantic reasoning and multi-turn editing; latent-diffusion iterations such as Imagen~4 scale the LDM recipe with stronger text encoders and curated data; preference-tuned systems such as Midjourney~v7~\cite{midjourney_v7_alpha} and FLUX.2~[max]~\cite{bfl_flux2_max} fine-tune strong priors on human-preference data (FLUX.2 on the rectified-flow recipe of §\ref{para:p3_s8_1_1_fm}; Midjourney's architecture is undisclosed).
Each reduces to one of the three paradigms plus some of the refinements just listed; because architectural and training detail is withheld, public benchmarking and reproduction still rest on the open SD / DiT stack.

\subsubsection{Benchmarks}
Three independent things are usually conflated under ``benchmark'': the training corpora on which models are fit, the evaluation suites of prompts that probe particular capabilities, and the metrics that score the resulting images.
A model trained on LAION-5B can be evaluated on GenEval prompts and scored by FID, CLIPScore, or an LLM judge, and disagreement between those scores is itself informative.

\paragraph{Training Corpora.}
Corpus scale has driven much of T2I progress.
Early work fit fine-grained sets with ten captions per image, CUB-200-2011 (11,788 bird images)~\cite{wah2011cub} and Oxford-102 Flowers (8,189 images)~\cite{nilsback2008flowers}; MS-COCO Captions~\cite{lin2014microsoft} (123,287 images, five captions each) became the standard general-scene corpus and, for years, the evaluation set as well, a conflation that the suites below were designed to end.
Current models train on web-scale image-text pairs: LAION-5B~\cite{schuhmann2022laion5b} releases 5.85 billion CLIP-filtered pairs, and proprietary corpora of similar scale, less documented, underlie DALL-E~2 and Parti.
Scale brought duplication, demographic bias, unsafe content, and licensing disputes with it, and filtered subsets such as LAION-Aesthetics and curation pipelines are now part of every training recipe.

\paragraph{Evaluation Benchmark Suites.}
An evaluation suite is a curated prompt set plus a verification protocol: prompts are adversarial or capability-stratified, reference images are usually absent, and scoring is automatic (an object detector, a VQA model, or an LLM/MLLM judge) rather than reconstruction-based.
Table~\ref{tab:p3_s8_1_1_eval_benchmarks} lists the main suites since 2022.
The earliest, DrawBench~\cite{saharia2022photorealistic} and PartiPrompts~\cite{parti2022}, are human-pairwise prompt sets released with Imagen and Parti; HRS-Bench~\cite{bakr2023hrsbench} and HEIM~\cite{lee2023heim} are holistic, scoring 13 and 12 axes respectively; TIFA~\cite{hu2023tifa}, GenEval~\cite{ghosh2023geneval}, and T2I-CompBench~\cite{huang2023t2icompbench} introduced automatic verification by VQA, object detection, and BLIP-VQA plus UniDet, and the latter two are now reported alongside FID and CLIPScore in most papers.
For spatial intelligence the salient suite is VISOR~\cite{gokhale2022visor}, which isolates two-object relations (left/right, above/below) over $25{,}280$ prompts and finds that otherwise strong models frequently render the wrong configuration.
DPG-Bench~\cite{hu2024ella} targets dense, long compositional prompts, GenAI-Bench~\cite{li2024genaibench} pairs compositional prompts with VQAScore judging, and WISE~\cite{niu2025wise} asks whether generated images agree with real-world facts (flag colors, landmarks, anatomy), a knowledge axis orthogonal to compositional binding.

\begin{table}[!ht]
\centering
\fontsize{9}{11}\selectfont
\setlength{\tabcolsep}{1.0mm}
\caption{Representative evaluation benchmark suites for text-to-image generation, ordered approximately by year of release. ``Scoring'' indicates the dominant verification mechanism: detector- or VQA-based automatic checking, learned preference models, MLLM/LLM-as-judge rubric scoring, or aggregated human preference. Several suites combine multiple scoring channels.}
\label{tab:p3_s8_1_1_eval_benchmarks}
\begin{tabular}{L{0.17\textwidth} C{0.06\textwidth} C{0.11\textwidth} L{0.33\textwidth} L{0.25\textwidth}}
\toprule
\textbf{Benchmark} & \textbf{Year} & \textbf{\#~Prompts} & \textbf{Evaluation Axis} & \textbf{Scoring} \\
\midrule
DrawBench~\cite{saharia2022photorealistic}       & 2022 & 200   & Broad capability + alignment ($11$ categories) & Human pairwise \\
PartiPrompts~\cite{parti2022}     & 2022 & 1{,}600 & Capability spectrum ($12$ categories) & Human pairwise \\
VISOR~\cite{gokhale2022visor}     & 2022 & 25{,}280 & Two-object spatial relations (left/right, above/below) & Object-detector verification \\
TIFA~\cite{hu2023tifa}            & 2023 & 4{,}081 & Faithfulness (VQA over auto-generated questions) & VQA / LLM-as-judge \\
HRS-Bench~\cite{bakr2023hrsbench} & 2023 & 45{,}000 & Holistic; $13$ skills $\times$ $50$ scenarios & Automatic, mixed \\
HEIM~\cite{lee2023heim}           & 2023 & 62 scenarios & Holistic; $12$ aspects (alignment, quality, knowledge, bias, toxicity, \ldots) & Automatic $+$ human \\
GenEval~\cite{ghosh2023geneval}   & 2023 & 553   & Object-centric: single/two-obj., counting, color, position, attribute & Detector verification \\
T2I-CompBench~\cite{huang2023t2icompbench} & 2023 & 6{,}000 & Compositional: attribute binding, object relationships, complex prompts & BLIP-VQA $+$ UniDet \\
DPG-Bench~\cite{hu2024ella}       & 2024 & 1{,}065 & Dense, long compositional prompts & MLLM-based scoring \\
GenAI-Bench~\cite{li2024genaibench} & 2024 & 1{,}600 & Compositional T2V/T2I prompts & VQAScore (LLM-judge) \\
WISE~\cite{niu2025wise}           & 2025 & 1{,}000 & World-knowledge-informed semantic accuracy & LLM-based grading \\
\bottomrule
\end{tabular}
\end{table}

\paragraph{Evaluation Metrics.}
Metrics fall into four tiers.
Distributional metrics compare the generated and real image populations and say nothing about any individual prompt: the Fr\'echet Inception Distance (FID)~\cite{heusel2017gans} is the Wasserstein-2 distance between Gaussian fits of Inception-V3 features of generated and real images, the Inception Score rewards images that are individually confidently classified yet collectively diverse, and generative precision/recall~\cite{sajjadi2018assessing,kynkaanniemi2019improved} separate quality from coverage through neighborhood tests in feature space.
Alignment metrics score one image against its prompt; CLIPScore~\cite{hessel2021clipscore}, the cosine similarity between the CLIP embeddings of $\hat{\mathbf{y}}$ and $\mathbf{x}$, is cheap and reference-free but inherits CLIP's insensitivity to attribute binding, counting, and spatial relations.
Learned preference models, HPSv2~\cite{wu2023hpsv2}, ImageReward~\cite{xu2023imagereward} (fit to $137$K human comparisons), and PickScore~\cite{kirstain2023pickapic} (fit to the ${\sim}500$K real-user comparisons of Pick-a-Pic), regress a scalar that correlates with average user preference, bundling realism, alignment, and aesthetics; they have become the default automatic ranking while inheriting their raters' biases and degrading out of distribution.

The newest tier uses a vision-language model as judge.
TIFA~\cite{hu2023tifa} has a language model write verification questions from the prompt and a VQA model answer them against the image, the score being the fraction answered correctly; VQAScore~\cite{lin2024evaluating} collapses this to a single yes/no probability under a fine-tuned CLIP-FlanT5 model and reports the best human correlation on compositional suites; VIEScore~\cite{ku2023viescore} and LLMScore~\cite{lu2023llmscore} prompt a multimodal LLM with a structured rubric and return per-axis scores with brief rationales, and direct prompting of frontier VLMs with custom rubrics has become routine for ad-hoc capability probes.
Judge metrics track human preference on compositional and spatial prompts far better than CLIPScore, at the cost of a frontier-model call per image, sensitivity to rubric wording, and drift across judge-model versions that makes numbers from different papers hard to compare.
No single tier suffices, and current papers report a tuple (FID or IS, CLIPScore, one or two preference scores, one judge score) on several suites plus a small human study; cross-metric disagreement is read as a warning rather than averaged away.


\subsubsection{Challenges and Trends}
The weakness that matters most for this handbook is \textbf{spatial-relation faithfulness}, a sub-problem of compositionality.
Models frequently misplace objects relative to one another: a prompt that puts a cup to the left of a laptop is often rendered with the two swapped or stacked, and VISOR~\cite{gokhale2022visor} shows that language-specified relations are frequently ignored outright.
Data-centric remedies help; SPRIGHT~\cite{chatterjee2024spright} re-captions training images with explicit spatial descriptions and measurably improves consistency, but reliable control over where objects appear remains open.
Because layout-conditioned generation (§\ref{subsec:p3_s8_1_2}) and text-to-3D (§\ref{sec:p3_s8_2}) inherit whatever relational competence this prior provides, the limitation is foundational rather than cosmetic.
Fine structure is a second gap: legible text, hands, and faces still fail in otherwise strong models and appear to need architectural or data-specific fixes rather than scale alone.

\textbf{Evaluation reliability} is the third.
Automatic scores correlate imperfectly with human judgment precisely on the compositional, counting, and spatial-relation prompts that matter here; judge-based metrics close part of the gap but add rubric sensitivity and judge-version drift, so headline numbers across papers are not directly comparable.
As GenEval and T2I-CompBench become standard, their prompts risk leaking into web-scraped training corpora, in which case high scores increasingly measure memorization rather than capability; periodic refresh of held-out prompts, blinded judges, and explicit train/test decontamination are the obvious countermeasures and are not yet routine.

\textbf{Societal concerns} follow the data.
Models inherit demographic bias from their corpora and will produce unsafe content on request; prompt classifiers, safety filters, and fine-tuning on curated subsets reduce but do not remove the problem, and photorealistic synthesis raises misinformation and intellectual-property questions that watermarking and data curation address only in part.
The extension of these priors along the time axis, where identity and motion must stay consistent across frames, is the subject of §\ref{sec:p3_s8_3}; the next two subsections stay in the image plane and ask how the prior can be made to obey spatial constraints and local edits.

